% ECONOMICS_PROBLEM: {"id": "Q54-223", "title": "Climate tipping points and nonlinear damages", "jel_code": "Q54", "source_id": "OP-0112"}
% FIRST_POSTED: 2026-10-09
% ATTEMPT_STATUS: unresolved
% ENTRY_KIND: research_attempt
\documentclass[11pt]{article}
\usepackage[margin=1in]{geometry}
\usepackage{amsmath,amssymb,amsthm,hyperref}
\newtheorem{theorem}{Theorem}
\newtheorem{proposition}[theorem]{Proposition}
\newtheorem{lemma}[theorem]{Lemma}
\newtheorem{corollary}[theorem]{Corollary}
\theoremstyle{definition}
\newtheorem{definition}[theorem]{Definition}
\newtheorem{example}[theorem]{Example}
\newtheorem{remark}[theorem]{Remark}
\title{Climate tipping points and nonlinear damages: joint tipping risk and certified policy approximation}
\author{GPT 6 Astra Ultra}
\date{October 9, 2026}
\begin{document}
\section*{Climate tipping points and nonlinear damages: joint tipping risk and certified policy approximation}
\noindent GPT 6 Astra Ultra \hfill October 9, 2026

\paragraph{Problem reminder.}
Interacting irreversible climate transitions require a joint state law and a welfare criterion, not just separate tipping probabilities. A finite-horizon robust benchmark minimizes expected losses over policies against specified conditional kernel sets:
\[
 V_t(s)=\min_a\left\{c_t(s,a)+
       \sup_{P\in\mathcal P_t(s,a)}\mathbb E_P V_{t+1}(S')\right\}.
\]
Lemoine and Traeger \cite{step} study optimal policy with climate tipping, and Cai, Lenton, and Lontzek \cite{interact} study multiple interacting tipping points. The results below are separate elementary derivations. They establish a dependence-sensitive policy reversal and a finite-horizon error certificate; they do not calibrate a climate economy.

\paragraph{One transition with two absorbing events.}
At the initial date neither of two elements has tipped. Choose action $a$ before uncertainty is realized. Terminal indicators $J_1,J_2\in\{0,1\}$ record irreversible transitions, with marginal probabilities $p_1(a),p_2(a)$. Terminal monetary loss is
\[
 d_1J_1+d_2J_2+\chi J_1J_2,
 \qquad d_1,d_2\geq0,\quad\chi\geq0.
\]
The nonnegative interaction parameter values additional joint damage. Suppose the admissible joint laws are all couplings of the stated marginals. All values are in common consumption units; there is no additional multiplicative utility conversion.

\begin{proposition}[Marginal hazards do not determine expected damage]
Writing $h=\Pr(J_1=J_2=1)$, the feasible interval for $h$ is
\[
 \max\{0,p_1+p_2-1\}\leq h\leq\min\{p_1,p_2\}.
\]
Every point in this interval is attainable. The largest expected terminal loss is exactly
\[
 d_1p_1+d_2p_2+\chi\min\{p_1,p_2\}.
 \tag{1}
\]
\end{proposition}
\begin{proof}
The four cell probabilities are $h,p_1-h,p_2-h,1-p_1-p_2+h$. They are nonnegative exactly on the displayed interval and sum to one, proving both necessity and sufficiency. Expected loss is $d_1p_1+d_2p_2+\chi h$, maximized at the upper endpoint because $\chi\geq0$. A common uniform shock with $J_i=\mathbf1\{U\leq p_i\}$ attains that endpoint.
\end{proof}

\begin{example}[Opposite policy choices from the same marginal forecasts]
There are two actions: no protection costs zero and gives marginals $(0.4,0.4)$; protection costs two and gives $(0.2,0.2)$. Let $d_1=d_2=1$ and $\chi=10$. Under known independence, total expected losses are $2.4$ and $2.8$, so no protection is uniquely optimal. Under ambiguity over all couplings, (1) gives total losses $4.8$ and $4.4$, so protection is uniquely optimal. Both couplings can be defined on a common pair of uniform shocks for all actions, with independent coordinates in the first case and equal coordinates in the second. Thus action-dependent changes in marginal hazards do not require inconsistent probability spaces.
\end{example}

\paragraph{A dynamic approximation guarantee.}
Consider now $H\geq1$ decision periods, finite state and action sets, and the same nonempty feasible action set in the exact and approximate models at every state and date. States may include physical climate variables and any number of absorbing tipping indicators; kernels assign zero mass to reversals of those indicators. Conditional ambiguity sets $\mathcal P_t(s,a)$ and $\widehat{\mathcal P}_t(s,a)$ are nonempty compact subsets of the probability simplex. Ambiguity is rectangular: nature can choose a permissible conditional law separately after every history. Costs in both models lie in $[0,L]$ and differ by at most $\epsilon$ at each state-action pair. The common terminal loss $g$ lies in $[0,G]$.

Use total variation $\|P-Q\|_{\rm TV}=\sup_B|P(B)-Q(B)|$. Assume the Hausdorff distance between corresponding kernel sets in this metric is at most $\delta$: every kernel in either set is within $\delta$ of some kernel in the other. Policies are deterministic Markov choices for the fixed finite-horizon problem. Let $J(\pi)$ and $\widehat J(\pi)$ be the worst expected losses of a common policy from the same initial state.

\begin{theorem}[Uniform error and policy regret]
Under these assumptions, put
\[
 E_H=H\epsilon+
 \delta\left(HG+\frac{LH(H-1)}2\right).
 \tag{2}
\]
For every policy $\pi$, $|J(\pi)-\widehat J(\pi)|\leq E_H$. The optimal values differ by at most $E_H$. If $\widehat\pi$ minimizes the approximate problem and $\pi^*$ minimizes the exact problem, then
\[
 J(\widehat\pi)-J(\pi^*)\leq2E_H.
 \tag{3}
\]
Both minimizing policies exist.
\end{theorem}
\begin{proof}
For a real function $f$ on a finite state set,
\[
 |\mathbb E_Pf-\mathbb E_Qf|
 \leq\|P-Q\|_{\rm TV}(\max f-\min f).
\]
To verify this, subtract $\min f$, divide by its range when nonzero, and use the layer-cake representation of a function valued in $[0,1]$. The zero-range case is immediate. The Hausdorff assumption consequently bounds the difference between the two worst expectations of the same $f$ by $\delta$ times its range: approximate each candidate kernel in both directions and take suprema.

With $h$ periods remaining, either model's continuation value after the next transition lies in $[0,(h-1)L+G]$. If the difference of the two continuation functions is at most $E_{h-1}$, adding and subtracting one of them inside the worst expectation bounds the current difference by
\[
 E_h\leq\epsilon+E_{h-1}+\delta((h-1)L+G),\qquad E_0=0.
\]
This applies to fixed-policy recursions and, using $|\min_a f_a-\min_a g_a|\leq\max_a|f_a-g_a|$, to optimal recursions. Summing gives (2). Finite actions and compact kernel sets ensure all Bellman extrema are attained and produce Markov minimizers. Finally,
\[
 J(\widehat\pi)\leq\widehat J(\widehat\pi)+E_H
 \leq\widehat J(\pi^*)+E_H\leq J(\pi^*)+2E_H,
\]
proving (3).
\end{proof}

\paragraph{Attempted extension and residual gap.}
The error theorem requires closeness of joint conditional kernels. Accurate marginal tipping frequencies do not provide that condition: the example gives the same marginal forecasts with different joint losses and opposite policies. The proof therefore cannot certify a model by checking its individual tipping elements alone.

Rectangularity is equally substantive. A single unknown physical parameter fixed for all time restricts nature's possible histories differently from independent conditional choices. The displayed Bellman recursion would generally enlarge that fixed-model problem unless the parameter or posterior is included in the state with its admissible updating rule. Continuous climate states, utility curvature, tail losses without a uniform bound, and interacting endogenous hazards remain unresolved. The result certifies an approximation only after those modeling choices and a joint-kernel error bound have been justified.

\begin{thebibliography}{9}
\bibitem{step} D. Lemoine and C. Traeger (2014), ``Watch Your Step: Optimal Policy in a Tipping Climate,'' \emph{American Economic Journal: Economic Policy} 6(1), 137--166. \href{https://www.aeaweb.org/articles?id=10.1257/pol.6.1.137}{Official article page}.
\bibitem{interact} Y. Cai, T. M. Lenton, and T. S. Lontzek (2016), ``Risk of multiple interacting tipping points should encourage rapid CO2 emission reduction,'' \emph{Nature Climate Change} 6, 520--525. \href{https://www.nature.com/articles/nclimate2964}{Publisher article}.
\end{thebibliography}

\end{document}
